
In this section we will study how the design parameters translate into
the actuation matrix $\mathbf{A}$ defined in the previous section. In
particular, we will consider how a bounded actuation translates into
maximum values of force and torque. We will then optimize these
parameters in order to maximize these values across all directions.

\subsection{Problem statement}
\label{sec:problem-statement}

% Throughout this paper we will assume that the propellers are equally
% distributed radially, that is,
% \begin{equation}
%   \label{eq:thetas}
%   \theta_i = (i-1) \frac{\pi}{3}
% \end{equation}
% For the sake of symmetry we will further consider the rotation
% directions to be alternating, that is,
% \begin{equation}
%   \label{eq:ws}
%   w_i =
%   \begin{cases}
%     1 & i\text{ is odd,} \\
%     -1 & i\text{ is even.}
%   \end{cases}
% \end{equation}

We will start by considering that each actuation signal is
bounded between $-1$ and $1$, that is,
\begin{equation}
  \label{eq:ahcube}
  -1 \leq u_i \leq 1 \quad \text{for} \quad i=1,\ldots,6
\end{equation}
According to \eqref{eq:au}, this hypercube will map onto a
6-dimensional convex polyhedron\footnote{A convex polyhedron is an
  intersection of a finite number of half-spaces~\cite{jeter86}.} in
the $(\bar{F}, \bar{M})$ space.  Any other choice of bounds is
possible by appropriately scaling constants $K_1$ and $K_2$.  However,
it assumes that the maximum propeller trust is symmetric with respect
to the direction of rotation. The remaining parameters are the angles
$\{\phi_i\}$. We will base our analysis on the optimization of these
angles with respect to various criteria.

Our goal will be to find the configurations of angles $\{\phi_i\}$
that maximize the range of forces (and torques) over all
directions. Conventional multirotors have all propellers directed at
the same direction, \textit{i.e.,} with $\{\phi_i=0\}$,, which
maximizes the thrust along that direction. However, they are unable to
produce any thrust component over an orthogonal direction. By having
nonzero values of $\{\phi_i\}$
we expect to gain thrust over any direction, at the
expense of reducing the maximum thrust over some particular one.
Geometrically, this corresponds to changing $\{\phi_i\}$
such that a ball of nonzero radius can fit inside the 3-dimensional
convex polyhedron in the $\bar{F}$ space mapped by the actuation
hypercube in~\eqref{eq:ahcube}, while keeping zero torque,
$\bar{M}=0$. A similar reasoning applies to the torque space
$\bar{M}$, while keeping $\bar{F}=0$. In sum, we will look for the values of $\{\phi_i\}$ that
maximize these radii.

First, we will address the problem of computing the maximum force
along a given direction specified as a unit vector $\hat{e}$, while
maintaining a zero torque.  From~\eqref{eq:au}, and assuming that
$\mathbf{A}$ is full rank, we get
\begin{equation}
  \bar{u} = \mathbf{A}^{-1}
  \left(
    \begin{array}{c}
      \bar{F} \\
      \bar{M}
    \end{array}
  \right) =\mathbf{A}^{-1}
  \left(
    \begin{array}{c}
      F \hat{e} \\
      0
    \end{array}
  \right)
\end{equation}
where $F>0$ is the force magnitude. For what follows, it will be
convenient to express the $\mathbf{A}^{-1}$ matrix as blocks of three
dimensional row vectors:
\begin{equation}
  \mathbf{A}^{-1} = \left[
    \begin{array}{cc}
      b_1^T & c_1^T \\
      \vdots & \vdots \\
      b_6^T & c_6^T \\
    \end{array}
    \right]
\end{equation}
where $b_i,c_i\in\mathbb{R}^3$. Then, the actuation of the $i$-th
propeller is given by
\begin{equation}
  u_i = F b_i^T \hat{e}
\end{equation}
Since $|u_i|\leq1$, we have $F|b_i^T\hat{e}|\leq1$, and thus $F$ has
this upper bound:
\begin{equation}
  F \leq \frac{1}{|b_i^T\hat{e}|}
\end{equation}
Since this inequality has to be satisfied for all propellers
$i=1,\ldots,6$, the maximum force $F^\mathrm{max}_{\hat{e}}$ is given
by the lowest of these upper bounds
\begin{equation}
  F^\mathrm{max}_{\hat{e}} = \min_i \frac{1}{|b_i^T\hat{e}|}
\end{equation}
This force is the maximum force along a given direction $\hat{e}$. The
maximum force attainable in any direction can be obtained by
minimising this force over all possible directions. Since
$|b_i^T\hat{e}|\leq\|b_i\|$, this minimum is given by
\begin{equation}
  \label{eq:fmax}
  F^\mathrm{max} = \min_i \frac{1}{\|b_i\|}
\end{equation}
The same reasoning can be applied to the torques: consider a torque
$\bar{M}=M\hat{e}$ along an arbitrary direction defined by $\hat{e}$,
the corresponding actuation with $\bar{F}=0$ is $u_i=M c_i^T \hat{e}$,
resulting in the following maximum torque along any direction:
\begin{equation}
  \label{eq:mmax}
  M^\mathrm{max} = \min_i \frac{1}{\|c_i\|}
\end{equation}

Now, these maximum forces and torque values depend on the design
parameters. In the following we will use an optimization approach to
find the values of these parameters that maximize the maximum force
and/or torque. We will consider the propellers to be equally
distributed radially, that is,
\begin{equation}
  \label{eq:thetas}
  \theta_i = (i-1) \frac{\pi}{3}
\end{equation}
and a fixed distance $d$, as well as the constants $K_1$ and $K_2$. All the
remaining parameters will be the unknown variables:
\begin{equation}
  \bar{\psi} = ( \phi_1, \ldots, \phi_6, w_1,\ldots, w_6 )^T
\end{equation}
with the feasibility domain defined by
\begin{equation}
  \Psi = \left\{ \bar{\psi} \::\: |\phi_{1,\ldots,6}| \leq
  \frac{\pi}{2}, w_{1,\ldots,6}\in\{-1,1\} \right\}
\end{equation}

Since we intend to both maximize force and torque, we will model this
problem as a multi-criteria optimization problem comprising two cost functions:
\begin{equation}
  \label{eq:multicriteria}
  J_1(\bar{\psi}) = -F^\mathrm{max} \qquad \mathrm{and} \qquad J_2(\bar{\psi}) = -M^\mathrm{max}
\end{equation}
using \eqref{eq:fmax} and \eqref{eq:mmax} computed from the inverse of
the $\mathbf{A}$ matrix defined in
section~\ref{sec:dynamical-model}. This matrix depends nonlinearly on
the unknowns $\bar{\psi}$. The solution of this multi-criteria
optimization problem is the set $P$ of non-dominated solutions:
$\bar{\psi}^0\in P$ if there is no $\bar{\psi}\in\Psi$ such that
$J_i(\bar{\psi})\leq J_i(\bar{\psi}^0)$ for all $i\in\{1,2\}$ and
$J_M(\bar{\psi})< J_M(\bar{\psi}^0)$ for at least one
$i\in\{1,2\}$. This set is also called \emph{Pareto optimal
  set}~\cite{statnikov95}. A method to approximate this set consists
in considering the linear combination of the~\eqref{eq:multicriteria}
cost functions, also known as \emph{linear scalarization}:
\begin{equation}
  \begin{aligned}
    \label{eq:linscal}
    J(\bar{\psi};\lambda) &= (1-\lambda) J_1(\bar{\psi}) + \lambda
    J_2(\bar{\psi}) \\
    &= -(1-\lambda) F^\mathrm{max} - \lambda M^\mathrm{max}
  \end{aligned}
\end{equation}
Since these problems cannot be solved in closed form, we will make use
of numerical optimization methods. The following section presents the
numerical results obtained for this problem.


% Considering the angles $\{\phi_i\}$ as the only design
% parameters, the following optimization problem can be formulated:
% \begin{equation}
%   \Phi^* = \min_{\Phi\in\mathcal{P}} J(\Phi)
% \end{equation}
% where $J$ is the cost function, $\Phi=[\phi_1\cdots\phi_6]^T$, and $\mathcal{P}$ is the set of
% allowable angle values. As cost function we consider a linear
% combination of the symmetric of the maximum force and torque,
% \begin{equation}
%   J(\Phi) = -(1-\lambda) F^\mathrm{max} - \lambda M^\mathrm{max}
% \end{equation}
% which corresponds to a linear scalarization of a multi-objective
% problem consisting of the objective functions~\eqref{eq:fmax}
% and~\eqref{eq:mmax}.  This problem is neither convex, smooth, not
% linear, so we will employ global optimization methods to solve it
% numerically. In particular, the results presented below were obtained
% using differential evolution algorithm~\cite{storn97}.

% Recall that the values of parameters $\{\theta_i\}$ and $\{w_i\}$ are
% set according to \eqref{eq:thetas} and \eqref{eq:ws}. For the
% $\{\phi_i\}$ we will consider three cases:
% \begin{enumerate}
% \item the simplest one consists in alternating a single angle value between
%   two symmetric values, similarly to what is done in~\cite{jiang13,voyles12}, that is
%   \begin{equation}
%     % \mathcal{P}_1 = \left\{ (\phi, -\phi, \phi, -\phi, \phi, -\phi)^T \::\:
%     %   \phi\in[0; \frac{\pi}{2}] \right\}
%     \mathcal{P}_1 = \left\{ \left(
%         \begin{array}{r}
%           \phi \\ -\phi \\ \phi \\ -\phi \\ \phi \\ -\phi
%         \end{array}
%       \right) \::\: |\phi|<\frac{\pi}{2}] \right\}
%   \end{equation}
% \item next we consider three degrees of freedom, corresponding to the
%   three principal axes joining pair os motors at opposite sides, that
%   is
%   \begin{equation}
%     % \mathcal{P}_3 = \left\{ (\phi_1, \phi_2, \phi_3, -\phi_1, -\phi_2, -\phi_3)^T \::\:
%     %   \phi_1\in[0; \frac{\pi}{2}], \phi_2,\phi_3\in[0; \frac{\pi}{2}] \right\}
%     \mathcal{P}_3 = \left\{ \left(
%         \begin{array}{r}
%           \phi_1 \\ \phi_2 \\ \phi_3 \\ -\phi_1 \\ -\phi_2 \\ -\phi_3
%         \end{array}
%       \right) \::\: |\phi_{1,\ldots,3}|<\frac{\pi}{2} \right\}
%   \end{equation}
% \item finally, there is the case where we have the six degrees of
%   freedom, one for each propeller, that
%   is
%   \begin{equation}
%     % \mathcal{P}_3 = \left\{ (\phi_1, \phi_2, \phi_3, -\phi_1, -\phi_2, -\phi_3)^T \::\:
%     %   \phi_1\in[0; \frac{\pi}{2}], \phi_2,\phi_3\in[0; \frac{\pi}{2}] \right\}
%     \mathcal{P}_3 = \left\{ \left(
%         \begin{array}{r}
%           \phi_1 \\ \phi_2 \\ \phi_3 \\ \phi_4 \\ \phi_5 \\ \phi_6
%         \end{array}
%       \right) \::\: |\phi_{1,\ldots,6}|<\frac{\pi}{2} \right\}
%   \end{equation}
% \end{enumerate}
% We have considered an unilateral interval for one of the propellers
% for symmetry breaking, since inverting the signs of all the $\phi_i$
% angles would result in the same cost function values.

\subsection{Numerical results}
\label{sec:numerical-results}

We applied this methodology to a vehicle with the design parameters
shown in Table~\ref{tab:params}. The constants $K_1$ and $K_2$ have
the same values than the ones reported in \cite{jiang13msc}, while $d$
is the one from the design presented in
Section~\ref{sec:specification}.

\begin{table}
  \centering
  \begin{tabular}{ccccc}
    $K_1$ & $K_2$ & $d$ & $\{\theta_i\}$  \\
    \hline
    5.7 & 1.3 & 0.16 & in \eqref{eq:thetas}  \\
  \end{tabular}
  \caption{Design parameters of the simulated vehicle.}
  \label{tab:params}
\end{table}

To get an idea of the shape of the feasible set in the space of costs
values, we sampled $N$ points uniformly from $\Psi$ and identified the
non-dominated points. The result is shown in
Figure~\ref{fig:pareto}. Although these non-dominated points are an
unstable approximate to the Pareto optimal set, they suggest that a
linear scalarization of the cost functions should provide a good
approximation to the actual Pareto optimal set.

\begin{figure}
  \centering
  \includegraphics[width=0.8\linewidth]{pareto_12.png}
  \caption{Samples ($N=10^7$) uniformly drawn from the feasible set $\Psi$ shown
    in the cost functions space. The non-dominated samples are shown
    as bigger points.}
  \label{fig:pareto}
\end{figure}

The shadow minimum is defined by the optimization of the individual
cost functions, that is,
\begin{equation}
  J_1^* = \min_{\bar{\psi}} J_1(\bar{\psi}) \qquad
  J_2^* = \min_{\bar{\psi}} J_2(\bar{\psi})
\end{equation}
After running differential evolution over~100 runs, we obtained the
minimum values $J_1^*=-11.40$ and $J_2^*=-3.176$ (with a total variation
below~$10^{-3}$).

% The shadow minimum is defined by the optimization of the individual
% cost functions, that is,
% \begin{equation}
%   \bar{\psi_1^*} = \argmin_{\bar{\psi}} J_1(\bar{\psi}) \qquad
%   \bar{\psi_2^*} = \argmin_{\bar{\psi}} J_2(\bar{\psi})
% \end{equation}
% These minimizers corresponds to two points in the costs function
% space, say $(J_1^*,J_2^+)=( J_1(\bar{\psi_1^*}),J_2(\bar{\psi_1^*}))$
% and $(J_1^+,J_2^*)=( J_1(\bar{\psi_2^*}),J_2(\bar{\psi_2^*}))$. The
% line between these two points is called the \emph{convex hull of
%   individual minima} (CHIM)}. If the Pareto optimal set is convex,
% which from Figure~\ref{fig:pareto} seems to be the case, then it
% certainly lies below the CHIM line~\cite{das98}. After running
% differential evolution over~100 runs, we obtained the mean values
% $(J_1^*,J_2^+)=(-11.40, -2.606)$ and
% $(J_1^+,J_2^*)=(-9.260 , -3.176)$.

We applied the linear scalarization~\eqref{eq:linscal} and used a
differential evolution\footnote{We used the implementation found in
  the \emph{optimize} package of SciPy. We used the default parameters
  except for a population size of 512 and a mutation dithering range
  of~$[0.2; 0.6]$.} algorithm~\cite{storn97} on this cost function,
while ranging $\lambda$ over an equally spaced grid from 0 to 1
totalling 20 runs. Figure~\ref{fig:pm12} shows the obtained
results. The bottom left solution in the \ref{fig:pm12}(b) plot, which
approximately dominates all of the others, includes 12 of these 20
runs. 

 \begin{figure}
  \centering
  % data from rm12dep_2.npy
  \subfloat[]{\includegraphics[width=0.33\linewidth]{out_pm12_cost.pdf}}
  \subfloat[]{\includegraphics[width=0.33\linewidth]{out_pm12_pareto.pdf}}
  \subfloat[]{\includegraphics[width=0.33\linewidth]{out_pm12_solution.pdf}}
  \caption{Optimization results,
    from left to right: (a)~plot of $J$, $J_1=-F^\mathrm{max}$, and
    $J_2=M^\mathrm{max}$ in function of $\lambda$, (b)~solutions found
    shown in the $J_1$ and $J_2$ space, and (c)~optimized $\{\phi_i\}$
    in function of~$\lambda$.}
  \label{fig:pm12}
\end{figure}

We then ran differential evolution 100 times with
$\lambda=0.4$, an heuristically chosen within a range of $\lambda$
values where the resulting costs seem stable. From these runs, the
dominant solution, found in 24 of these 100 runs, is shown in
Table~\ref{tab:best}. Its costs are $J_1=-11.40$ and $J_2=-3.176$,
which coincide with the individual minima $J_1^*$ and $J_2^*$ found
above. That is, no other solution found has lower cost in either of
the cost functions. 

\begin{table}
  \centering
  \begin{tabular}{ccccccc}
    propeller ($i$) & 1 & 2 & 3 & 4 & 5 & 6 \\
    \hline
    $\theta_i$ & 0\textdegree & 60\textdegree & 120\textdegree & 180\textdegree & 240\textdegree & 300\textdegree \\
    $\phi_i$ & -74.54\textdegree & -39.54\textdegree & -74.57\textdegree & -39.48\textdegree & -74.49 \textdegree & -39.46 \\
    $w_i$ & -1 & 1 & -1 & 1 & -1 & 1 \\
    \hline
  \end{tabular}
  \caption{The solution found over 100 runs of the optimizer with
    $\lambda=0.4$.}
  \label{tab:best}
\end{table}

To wrap up these results, Figure~\ref{fig:costs} shows the solutions
found by differential evolution as boxplots, in three different
scenarios: the solutions for the two shadow minima
$\bar{\psi}_1^*=\arg\min J_1(\bar{\psi})$ and
$\bar{\psi}_2^*=\arg\min J_2(\bar{\psi})$, and the solutions
$\bar{\psi}_\lambda^*=\arg\min J(\bar{\psi}; \lambda)$ for
$\lambda=0.4$. Note that the minimum costs of the
$\bar{\psi}_\lambda^*$ solutions coincide with the individual
minima of $J_1^*=J_1(\bar{\psi}_1^*)$ and $J_2^*=J_2(\bar{\psi}_2^*)$.

\begin{figure}
  \centering
  \includegraphics[width=0.8\linewidth]{out_osom12.pdf}
  \caption{Boxplots of the $J_1$ and $J_2$ costs found while
    separately minimizing, $J_1(\bar{\psi})$ (first and second
    boxplot), $J_2(\bar{\psi})$ (third and fourth boxplot), and
    $J(\bar{\psi}; \lambda)$ for $\lambda=0.4$ (fifth and
    sixth). Each boxplots represents the minimum, maximum, and median
    values.}
  \label{fig:costs}
\end{figure}

The solution in Table~\ref{tab:best} implies three of the propellers to be
almost horizontal, hindering the air flow to produce the desired
thrust. Therefore, we decided on a different solution, with the angle
values rounded from an handpicked solution from the Pareto optimal set
obtained above. Its parameters are shown in Table~\ref{tab:selected}
and its costs are $J_1=-F^\mathrm{max}=-11.40$ and
$J_2=-M^\mathrm{max}=-2.623$, that is, the same maximum force but with
82\% of the maximum torque. We preferred this solution
because the propeller angles are not as horizontal, resulting in a
significantly better air flow.
% Running the
% optimization algorithm on a restricted feasibility space parametrized
% by~$\phi$ defined as
% \begin{equation}
%   \bar{\psi} = \left(-\phi, \phi, -\phi, \phi, -\phi, \phi, -1, -1, 1, 1,
%     1, -1\right)^T
% \end{equation}
% with $\lambda=0.4$ resulted 

\begin{table}
  \centering
  \begin{tabular}{ccccccc}
    propeller ($i$) & 1 & 2 & 3 & 4 & 5 & 6 \\
    \hline
    $\theta_i$ & 0\textdegree & 60\textdegree & 120\textdegree & 180\textdegree & 240\textdegree & 300\textdegree \\
    $\phi_i$ & -55\textdegree & 55\textdegree & -55\textdegree & 55\textdegree & -55\textdegree & 55\textdegree \\
    $w_i$ & -1 & -1 & 1 & 1 & 1 & -1 \\
    \hline
  \end{tabular}
  \caption{Design parameters of the selected solution.}
  \label{tab:selected}
\end{table}


% Figure~\ref{fig:forces} shows the maximum forces and torques in set
% $\mathcal{P}_1$. The cost functions~\eqref{eq:fmax}
% and~\eqref{eq:mmax} correspond to the symmetric of the pointwise
% minimum over the three lines for each one of these plots. The
% non-convexity of the cost functions is clearly visible in this plot. When
% $\phi=0$ the maximum force is zero, since in this configuration all
% thrust is directed along the $Z$ axis: this corresponds to the
% usual configuration of hexacopters (and quadrotors, although with less
% propellers). As $|\phi|$ increases, the maximum force along directions
% other than $Z$ increases.

% \begin{figure}
%   \centering
%   \includegraphics[width=\linewidth]{out_plot_forces.pdf}
%   \caption{Maximum force and torque, $F^\mathrm{max}$ and
%     $M^\mathrm{max}$, for each one of
%     the body frame axes while ranging across the design space
%     $\mathcal{P}_1$.}
%   \label{fig:forces}
% \end{figure}


% \begin{figure}
%   \centering
%   \includegraphics[width=0.32\linewidth]{out_pm1_cost.pdf} \hfill
%   \includegraphics[width=0.32\linewidth]{out_pm1_curve.pdf} \hfill
%   \includegraphics[width=0.32\linewidth]{out_pm1_solution.pdf} 
%   \caption{Optimization results for feasibility set $\mathcal{P}_1$,
%     from left to right: (a)~plot of $-J$, $F^\mathrm{max}$, and
%     $M^\mathrm{max}$ in function of $\lambda$, (b)~relationship between $F^\mathrm{max}$ and
%     $M^\mathrm{max}$, and (c)~optimized $\phi$ in function of $\lambda$.}
%   \label{fig:pm1}
% \end{figure}

% \begin{figure}
%   \centering
%   \includegraphics[width=0.32\linewidth]{out_pm3_cost.pdf} \hfill
%   \includegraphics[width=0.32\linewidth]{out_pm3_curve.pdf} \hfill
%   \includegraphics[width=0.32\linewidth]{out_pm3_solution.pdf} 
%   \caption{Optimization results for feasibility set $\mathcal{P}_3$,
%     same description as in Figure~\ref{fig:pm1}}
%   \label{fig:pm3}
% \end{figure}

%%% Local Variables:
%%% mode: latex
%%% TeX-master: "main"
%%% End:
